\clearpage
\section*{2026-11-10 复核更正：任务一当前模型}
\markboth{任务一复核更正}{任务一复核更正}
\addcontentsline{toc}{section}{复核更正：任务一当前模型}
\textbf{适用范围。}本节替代前文任务一的水分定义、标定口径及其向后续任务的接口。
任务二现按2026-11-11更正节作情景实现；任务三仍为待验证方案，不能视为已经求解。本文只规定模型、参数和约束，不规定编程交付物。
题目并未给定数值权重、灌溉定额或工程预算；以下情景值均不冒充题干常数。

\paragraph{两个状态与尺度。}
$V_y$ 是年度生长季（6--9 月）NDVI 派生覆盖指标，沿用端元 $0.05/0.85$，
不是直接测量的地面覆盖率或年末覆盖率。$M_j$ 是根区高于凋萎点的可利用储水量（mm），
$0\le M_j\le C$；$M_j/C$ 是有效水分比例。年内固定 $V_y$ 计算水分过程，年末更新植被：
\[
 V_{y+1}=V_y+rV_y(1-V_y/K)\overline W_y-dV_y
             -kE_y(1-\eta V_y/K)V_y .
\]
年尺度温度影响并入有效增长率 $r$，不得再将年平均温度因子乘入生长项。
温度响应仍用于潜在蒸腾的季节分配；这是一种简化假设，不是温度恒定的实证结论。

\paragraph{储水闭合。}
每个月的降水、参考蒸散及分配到该月的灌溉均匀分配到当月各天，作为数值子步假设，
不能称为真实逐日天气。子步 $j$ 的潜在蒸腾为
\[
 D_j=c\min(1,2f_{c,y})f(T_j)ET_{0,j},\qquad
 f_{c,y}=\operatorname{clip}(V_y/V_{\rm full},0,1).
\]
$f_c$ 只是采用局地端元的覆盖代理；遥感绿洲高分位不证明真实冠层完全闭合。
设毛灌溉 $I_j$、有效降水比例 $\alpha_P$、净可蒸腾灌溉比例 $\eta_w$、地下补给 $G_j$，则
\begin{align*}
 B_j&=M_j+\alpha_PP_j+\eta_wI_j+G_j,\\
 Q_j&=\max(B_j-C,0),\qquad B'_j=B_j-Q_j,\\
 K_{s,j}&=\min\{1,B'_j/[(1-p)C]\},\qquad
 T_{a,j}=\min(B'_j,K_{s,j}D_j),\\
 M_{j+1}&=B'_j-T_{a,j},\qquad
 \overline W_y=\frac{\sum_{j\in y}T_{a,j}}{\sum_{j\in y}D_j}.
\end{align*}
分母为零时取 $\overline W_y=0$。每一步严格满足
$M_{j+1}-M_j=\alpha_PP_j+\eta_wI_j+G_j-T_{a,j}-Q_j$。
截留、土壤蒸发和无效灌溉损失已包含在有效入水比例中，不再重复扣除。
储水与胁迫的结构参考 FAO-56 第 8 章；其具体离散化和系数是本队情景设定。

\paragraph{参数的证据等级。}
$r,d$ 在训练目标年不晚于 2022 年的站点转移上合并拟合；$K$ 仍按各站训练期
$1.15Q_{0.95}(V)$ 确定。灌溉绿洲标定时假定 $W=1$，尚无灌溉实测验证这一假定。
$k=0,\eta=1$ 是基准情景，$k$ 的情景范围为 $0$--$0.10$；不能把零估计解释为风蚀不存在。
水分基准假设为 $C=75$ mm（范围 $30$--$150$）、$p=0.5$（$0.3$--$0.7$）、
$c=1$（$0.7$--$1.2$）、$\eta_w=0.75$（$0.65$--$0.85$）、
$\alpha_P=0.25$（$0$--$1$）、$G=0$。这些值没有由不完整的 SoilGrids 数据拟合。
\textbf{不得宣称水分响应已获现场标定。}

\paragraph{数据约束与可识别性。}
MODIS 必须逐像元匹配 QA，仅保留可靠性为 0 或 1 的像元；合成期有效像元不少于一半，
并用 75\% 门槛作敏感性检查。年度覆盖需至少两个有效生长季月份；插补值不用于标定和评分。
站点准入按训练期生长季缺失比例不超过 20\% 判断，不用验证期成绩筛站。
风指数先对日最大风速计算非线性代理 $q\propto u^2\max(u-u_T,0)$，再聚合，
以各站训练期日代理均值归一；它不是实测输沙通量或沙尘暴发生次数。

在不受水限制且 $k=0$ 的区段自由拟合中，递推式仅含 $a=r-d$ 和 $b=r/K$，
所以自由的 $r,d,K$ 不可分别识别。区段像元按后期变绿筛选，建立年份也使用未来观测，
因此 CON4 只能作回溯空间诊断，不能作为独立预测验证。

\paragraph{稳定性、区段和后续接口。}
年步长的安全计算域为 $0\le r\le1$、$d+kE\le1$、$0\le\eta\le1$、$0\le V\le K$。
域外参数应缩短时间步长并重新标定，不能截断掩盖问题。
无水且 $k=0$ 时的离散半衰期为 $\log(1/2)/\log(1-d)$，不是连续时间的 $\log2/d$。
给定 $W$ 的旧解析稳态式只是条件关系；水分与覆盖相互依赖时须求耦合周期不动点并检查收敛。
后续任务必须传递 $(V,M)$ 状态，不能将新植区立即赋予稳态覆盖，也不能仅用静态曲线代替状态转移。
研究对象为于田--民丰南缘的既有选窗沿线；1 km 宽度仅是代表性研究情景，
不代表已核实的工程边界或可治理面积。水量按 mm 和每公顷输出，$1$ mm $=10$ m$^3$/ha，
不得在缺乏真实区段面积时推断全工程用水量。
